% Source: https://artofproblemsolving.com/community/c3h2188094_some_fun_problemweek_1
% Stable identity: aops-topic-2188094-q05
Let $d$ be the length of a diagonal. The interior angles are $108^\circ$, and the diagonals from a vertex split its angle into three angles of $36^\circ$ (each is an inscribed angle on an arc of $72^\circ$ of the circumcircle). Let the diagonal $AC$ meet $BE$ at $P$ and $BD$ at $Q$.

\begin{asy}
usepackage("amsmath");
size(5cm);
defaultpen(fontsize(10pt));
pair[] V; for(int k=0;k<5;++k) V[k]=dir(90+72*k);
pair[] I; for(int k=0;k<5;++k) I[k]=extension(V[k],V[(k+2)%5],V[(k+1)%5],V[(k+4)%5]);
fill(I[0]--I[1]--I[2]--I[3]--I[4]--cycle,paleyellow);
draw(V[0]--V[1]--V[2]--V[3]--V[4]--cycle,linewidth(1));
for(int k=0;k<5;++k) draw(V[k]--V[(k+2)%5],gray);
draw(I[0]--I[1]--I[2]--I[3]--I[4]--cycle,red+linewidth(1));
label("$A$",V[0],N); label("$B$",V[1],W); label("$C$",V[2],SW); label("$D$",V[3],SE); label("$E$",V[4],E);
pair P=extension(V[1],V[4],V[0],V[2]), Q=extension(V[1],V[3],V[0],V[2]);
dot(P); dot(Q); label("$P$",P,NE,fontsize(8pt)); label("$Q$",Q,dir(10),fontsize(8pt));
\end{asy}

\textbf{The diagonal.} In triangle $ABP$, $\angle BAP=\angle ABP=36^\circ$, so it is similar to triangle $ACB$ (angles $36^\circ,36^\circ,108^\circ$). Hence $\dfrac{AP}{AB}=\dfrac{AB}{AC}$, so $AP=\dfrac{n^2}{d}$. In triangle $ABQ$, $\angle BAQ=36^\circ$ and $\angle ABQ=72^\circ$, so $\angle AQB=72^\circ$ and $AQ=AB=n$; by symmetry $CP=n$ as well. So $AP=d-n$, and $(d-n)d=n^2$, giving $d=\dfrac{1+\sqrt5}{2}\,n$.

\textbf{The inner side.} The inner side is $PQ=AQ-AP=n-(d-n)=2n-d=\dfrac{3-\sqrt5}{2}\,n$.

\textbf{The area.} The inner pentagon is regular, so it is similar to $ABCDE$ with ratio $\dfrac{3-\sqrt5}{2}$, and its area is $\left(\dfrac{3-\sqrt5}{2}\right)^2=\dfrac{7-3\sqrt5}{2}$ times the area of $ABCDE$. A regular pentagon with side $s$ has area $\dfrac{s^2}{4}\sqrt{25+10\sqrt5}$ (five triangles from the centre, each with base $s$ and height $\dfrac{s}{2}\cot36^\circ$). Therefore the inner pentagon has area
\[\dfrac{7-3\sqrt5}{2}\cdot\dfrac{n^2}{4}\sqrt{25+10\sqrt5}=\boxed{\dfrac{(7-3\sqrt5)\sqrt{25+10\sqrt5}}{8}\,n^2}=\dfrac{\sqrt{250-110\sqrt5}}{8}\,n^2\approx0.251\,n^2 .\]
